% GENERATED FILE. Edit canonical lessons, then run npm run generate:books.
% canonical-source-hash: fnv1a64:48552deb38a9c02d
% canonical-lessons: SA-C416-navikah, SA-C416-aritram, SA-C416-karnadharah, SA-C416-vahayati, SA-C416-uttarati

\chapter{Boatmen}
\label{ch:sa-boatmen}
\hlchaptermodality{416}

\begin{chapteropening}
\textbf{By the end of this chapter:} \emph{I can say who steers and rows the boat and where we land.}

Complete the last lesson of chapter 416: I can say who steers and rows the boat and where we land.
\end{chapteropening}

\section[\texorpdfstring{\sk{नाविकः} (nāvikaḥ)}{नाविकः (nāvikaḥ)}]{\sk{नाविकः} (nāvikaḥ) — a boatman}

\label{lesson:SA-C416-navikah}



\begin{quote}
Before the new one: say the Sanskrit for the far bank, then the Sanskrit for a causeway.
\end{quote}

\subsection*{You'll want to know: \sk{नाविकः}}
\textbf{\sk{नाविकः}} — \emph{nāvikaḥ} — \textquotedblleft{}a boatman\textquotedblright{}.

\begin{quote}
\sk{नाविकः} \sk{नौकां} \sk{नदीतीरम्} \sk{आनयति।}
\end{quote}

\begin{quote}
\emph{nāvikaḥ naukāṃ nadītīram ānayati.}
\end{quote}

\begin{quote}
The boatman brings the boat to the riverbank.
\end{quote}

A masculine noun.

\begin{cousinweb}
\textbf{\sk{नाविकः}} (\emph{nāvikaḥ}) is built on \emph{nau}, \textquotedblleft{}a boat\textquotedblright{}, with the ending \emph{-ika}; MW gives \emph{nāvika} as a helmsman, a pilot, a sailor (MW).
\end{cousinweb}

\begin{grammarlens}[title={What you've built}]
One more word for getting around.
\end{grammarlens}

\subsection*{Your turn}
\begin{itemize}
\raggedright
  \item \emph{Say it:} \emph{nāvikaḥ}
  \item \emph{Say it:} \emph{nāvikaḥ}, once more
  \item \emph{Say it:} \emph{nāvikaḥ naukāṃ nadītīram ānayati.} — The boatman brings the boat to the riverbank.
  \item \emph{Recall:} say the Sanskrit for the far bank, then the Sanskrit for a causeway, then say \emph{nāvikaḥ} again
\end{itemize}

\begin{tcolorbox}[breakable,colback=teal!4,colframe=teal!35!black,title={Before you move on}]
What does \sk{नाविकः} mean? (\textquotedblleft{}A boatman\textquotedblright{}.) Say it once more.
\end{tcolorbox}
\section[\texorpdfstring{\sk{अरित्रम्} (aritram)}{अरित्रम् (aritram)}]{\sk{अरित्रम्} (aritram) — an oar}

\label{lesson:SA-C416-aritram}



\begin{quote}
Before the new one: say the Sanskrit for a causeway, then the Sanskrit for a boatman.
\end{quote}

\subsection*{You'll want to know: \sk{अरित्रम्}}
\textbf{\sk{अरित्रम्}} — \emph{aritram} — \textquotedblleft{}an oar\textquotedblright{}.

\begin{quote}
\sk{नाविकस्य} \sk{हस्ते} \sk{अरित्रम्} \sk{अस्ति।}
\end{quote}

\begin{quote}
\emph{nāvikasya haste aritram asti.}
\end{quote}

\begin{quote}
The boatman has an oar in his hand.
\end{quote}

A neuter noun.

\begin{cousinweb}
\textbf{\sk{अरित्रम्}} (\emph{aritram}) is built on the root √ṛ, \textquotedblleft{}to move\textquotedblright{}, with the ending \emph{-itra}: MW gives \emph{aritra} as propelling and as an oar (MW).
\end{cousinweb}

\begin{grammarlens}[title={What you've built}]
One more word for getting around.
\end{grammarlens}

\subsection*{Your turn}
\begin{itemize}
\raggedright
  \item \emph{Say it:} \emph{aritram}
  \item \emph{Say it:} \emph{aritram}, once more
  \item \emph{Say it:} \emph{nāvikasya haste aritram asti.} — The boatman has an oar in his hand.
  \item \emph{Recall:} say the Sanskrit for a causeway, then the Sanskrit for a boatman, then say \emph{aritram} again
\end{itemize}

\begin{tcolorbox}[breakable,colback=teal!4,colframe=teal!35!black,title={Before you move on}]
What does \sk{अरित्रम्} mean? (\textquotedblleft{}An oar\textquotedblright{}.) Say it once more.
\end{tcolorbox}
\section[\texorpdfstring{\sk{कर्णधारः} (karṇadhāraḥ)}{कर्णधारः (karṇadhāraḥ)}]{\sk{कर्णधारः} (karṇadhāraḥ) — a helmsman}

\label{lesson:SA-C416-karnadharah}



\begin{quote}
Before the new one: say the Sanskrit for a boatman, then the Sanskrit for an oar.
\end{quote}

\subsection*{You'll want to know: \sk{कर्णधारः}}
\textbf{\sk{कर्णधारः}} — \emph{karṇadhāraḥ} — \textquotedblleft{}a helmsman\textquotedblright{}.

\begin{quote}
\sk{कर्णधारः} \sk{नौकां} \sk{सम्यक्} \sk{नयति।}
\end{quote}

\begin{quote}
\emph{karṇadhāraḥ naukāṃ samyak nayati.}
\end{quote}

\begin{quote}
The helmsman steers the boat well.
\end{quote}

A masculine noun.

\begin{cousinweb}
\textbf{\sk{कर्णधारः}} (\emph{karṇadhāraḥ}) is built from \emph{karṇa}, \textquotedblleft{}an ear\textquotedblright{}, which MW also gives as the helm or rudder of a ship, and \emph{dhāra}, \textquotedblleft{}holding\textquotedblright{}; MW gives the compound as a helmsman, a pilot (MW).
\end{cousinweb}

\begin{grammarlens}[title={What you've built}]
One more word for getting around.
\end{grammarlens}

\subsection*{Your turn}
\begin{itemize}
\raggedright
  \item \emph{Say it:} \emph{karṇadhāraḥ}
  \item \emph{Say it:} \emph{karṇadhāraḥ}, once more
  \item \emph{Say it:} \emph{karṇadhāraḥ naukāṃ samyak nayati.} — The helmsman steers the boat well.
  \item \emph{Recall:} say the Sanskrit for a boatman, then the Sanskrit for an oar, then say \emph{karṇadhāraḥ} again
\end{itemize}

\begin{tcolorbox}[breakable,colback=teal!4,colframe=teal!35!black,title={Before you move on}]
What does \sk{कर्णधारः} mean? (\textquotedblleft{}A helmsman\textquotedblright{}.) Say it once more.
\end{tcolorbox}
\section[\texorpdfstring{\sk{वाहयति} (vāhayati)}{वाहयति (vāhayati)}]{\sk{वाहयति} (vāhayati) — rows (a boat); drives}

\label{lesson:SA-C416-vahayati}



\begin{quote}
Before the new one: say the Sanskrit for an oar, then the Sanskrit for a helmsman.
\end{quote}

\subsection*{You'll want to know: \sk{वाहयति}}
\textbf{\sk{वाहयति}} — \emph{vāhayati} — \textquotedblleft{}rows (a boat); drives\textquotedblright{}.

\begin{quote}
\sk{नाविकः} \sk{अरित्रेण} \sk{नौकां} \sk{वाहयति।}
\end{quote}

\begin{quote}
\emph{nāvikaḥ aritreṇa naukāṃ vāhayati.}
\end{quote}

\begin{quote}
The boatman rows the boat with an oar.
\end{quote}

Past (he …): \sk{सः} \sk{नौकां} \sk{वाहितवान्}, \emph{saḥ naukāṃ vāhitavān}.

\begin{cousinweb}
\textbf{\sk{वाहयति}} (\emph{vāhayati}) is built on the root √vah, \textquotedblleft{}to carry\textquotedblright{}, which gives \emph{vahati}, \textquotedblleft{}carries\textquotedblright{}; MW gives its causative as to make carry, to drive a chariot and to propel a boat, and KRM gives the participle \emph{vāhita} (MW).
\end{cousinweb}

\begin{grammarlens}[title={What you've built}]
One more word for getting around.
\end{grammarlens}

\subsection*{Your turn}
\begin{itemize}
\raggedright
  \item \emph{Say it:} \emph{vāhayati}
  \item \emph{Say it:} \emph{vāhayati}, once more
  \item \emph{Say it:} \emph{nāvikaḥ aritreṇa naukāṃ vāhayati.} — The boatman rows the boat with an oar.
  \item \emph{Recall:} say the Sanskrit for an oar, then the Sanskrit for a helmsman, then say \emph{vāhayati} again
\end{itemize}

\begin{tcolorbox}[breakable,colback=teal!4,colframe=teal!35!black,title={Before you move on}]
What does \sk{वाहयति} mean? (\textquotedblleft{}Rows (a boat); drives\textquotedblright{}.) Say it once more.
\end{tcolorbox}
\section[\texorpdfstring{\sk{उत्तरति} (uttarati)}{उत्तरति (uttarati)}]{\sk{उत्तरति} (uttarati) — comes ashore; lands}

\label{lesson:SA-C416-uttarati}



\begin{quote}
Before the new one: say the Sanskrit for a helmsman, then the Sanskrit for rows.
\end{quote}

\subsection*{You'll want to know: \sk{उत्तरति}}
\textbf{\sk{उत्तरति}} — \emph{uttarati} — \textquotedblleft{}comes ashore; lands\textquotedblright{}.

\begin{quote}
\sk{यात्रिकाः} \sk{नौकायाः} \sk{तीरे} \sk{उत्तरन्ति।}
\end{quote}

\begin{quote}
\emph{yātrikāḥ naukāyāḥ tīre uttaranti.}
\end{quote}

\begin{quote}
The travellers come ashore from the boat.
\end{quote}

Past (he …): \sk{सः} \sk{नौकायाः} \sk{उत्तीर्णवान्}, \emph{saḥ naukāyāḥ uttīrṇavān}.

\begin{cousinweb}
\textbf{\sk{उत्तरति}} (\emph{uttarati}) is built from \emph{ud}, \textquotedblleft{}up, out\textquotedblright{}, and the root √tṝ, \textquotedblleft{}to cross\textquotedblright{}; MW gives \emph{ut-tṝ} as to pass out of the water, to disembark, with the participle \emph{uttīrṇa}, \textquotedblleft{}landed\textquotedblright{} (MW).
\end{cousinweb}

\begin{grammarlens}[title={What you've built}]
One more word for getting around.
\end{grammarlens}

\subsection*{Your turn}
\begin{itemize}
\raggedright
  \item \emph{Say it:} \emph{uttarati}
  \item \emph{Say it:} \emph{uttarati}, once more
  \item \emph{Say it:} \emph{yātrikāḥ naukāyāḥ tīre uttaranti.} — The travellers come ashore from the boat.
  \item \emph{Recall:} say the Sanskrit for a helmsman, then the Sanskrit for rows, then say \emph{uttarati} again
\end{itemize}

\begin{tcolorbox}[breakable,colback=teal!4,colframe=teal!35!black,title={Before you move on}]
What does \sk{उत्तरति} mean? (\textquotedblleft{}Comes ashore; lands\textquotedblright{}.) Say it once more.
\end{tcolorbox}
